% ============ §29 The structure theorems for complete local rings (原书 pp.223-230) ============
\section{The structure theorems for complete local rings}\label{sec:29}\index{coefficient ring, section 29}\index{p-ring, section 29}\index{quasi-coefficient ring, section 29}\index{Eisenstein extension, section 29}\index{ramified local ring, section 29}\index{unramified local ring, section 29}

By Theorem~\hyperref[thm:28.3]{28.3}, a complete local ring of equal characteristic $A$ has a
coefficient field. If $K$ is a coefficient field of $A$, and $x_1, \ldots, x_n$ are generators
of the maximal ideal, then any element of $A$ can be expanded as a power series
in $x_1, \ldots, x_n$ with coefficients in $K$, and therefore $A$ is a quotient of the
regular local ring $K[[X_1, \ldots, X_n]]$. We now want to extend all of this to the
case of unequal characteristic.

We say that a DVR of characteristic $0$ is a \emph{$p$-ring}\index{p-ring} if its maximal ideal is
generated by the prime number $p$. If $K$ is a given field of characteristic $p$, then
there exists a $p$-ring having $K$ as its residue field. This follows by applying the
next theorem to $A = \mathbb{Z}_{(p)}$.

\begin{theorem}\label{thm:29.1}
Let $(A, tA, k)$ be a DVR and $K$ an extension field of $k$; then there exists a DVR
$(B, tB, K)$ containing $A$.
\end{theorem}

\begin{proof}
Let $\{x_\lambda\}_{\lambda \in \Lambda}$ be a transcendence basis of $K$ over $k$, and set
$k_1 = k(\{x_\lambda\})$. We take indeterminates $\{X_\lambda\}_{\lambda \in \Lambda}$ over $A$ in
bijection with the $\{x_\lambda\}$, and set $A[\{X_\lambda\}] = A'$ and $A_1 = (A')_{tA'}$. Now
$A'$ is a free $A$-module, and hence separated for the $t$-adic topology; therefore,
so is $A_1$, and $A_1$ is a DVR with $A_1/tA_1 \cong k_1$. Hence, replacing $A$ and $k$ by $A_1$
and $k_1$, we can assume that $K$ is algebraic over $k$. Let $L$ be an algebraic
closure of the field of fractions of $A$, and let $\mathcal{F}$ be the set of all pairs
$(B, \varphi)$, where $B$ is an intermediate ring $A \subset B \subset L$ and
$\varphi\colon B \to K$ is an $A$-algebra homomorphism, satisfying the conditions
\[
  (*) \quad B \text{ is a DVR and } \Ker \varphi = \rad(B) \text{ is generated by } t.
\]
We introduce an order on $\mathcal{F}$ by defining $(B, \varphi) \le (C, \psi)$ if $B \subset C$ and
$\psi|_B = \varphi$. Let us show that $\mathcal{F}$ has a maximal element. Suppose that
$\{(B_i, \varphi_i)\}$ is a totally ordered subset of $\mathcal{F}$, and set
\[
  B_0 = \bigcup_i B_i, \quad \varphi_0 = \bigcup_i \varphi_i;
\]
then one sees easily that $B_0$ is a local ring with maximal ideal $tB_0$. If
$0 \ne x \in B_0$, then $x \in B_i$ for some $i$, and since $B_i$ is a DVR we can write
$x = t^\nu u$ with $u$ a unit of $B_i$ and some $\nu$. From this we get
$x \notin t^{\nu+1}B_0$, so that $B_0$ is $t$-adically separated. Hence $B_0$ is a DVR, and if
$\varphi_0\colon B_0 \to K$ is defined to be equal to $\varphi_i$ on $B_i$ then
$(B_0, \varphi_0) \in \mathcal{F}$. Hence by Zorn's lemma $\mathcal{F}$ has a maximal element;
suppose that $(B, \varphi)$ is one. Then if $\varphi(B) \ne K$, take $\alpha \in K - \varphi(B)$,
let $\bar{f}(X)$ be the minimal polynomial of $\alpha$ over $\varphi(B)$, and take a manic
$f(X) \in B[X]$ which is an inverse image of $\bar{f}(X)$. Then $f(X)$ is irreducible in
$B[X]$, and hence (by Ex.~\ref{ex:9.6-ex}) also irreducible over the field of fractions of
$B$. Let $\xi$ be a root of $f(X)$ in $L$, and set $B' = B[\xi]$; then
$B' = B[X]/(f)$. Therefore
\[
  B'/tB' = B[X]/(t, f) = \varphi(B)[X]/(\bar{f}) = \varphi(B)(\alpha)
\]
is a field; since $B'$ is integral over $B$, every maximal ideal of $B'$ contains
$tB'$, so that $B'$ is a local integral domain with maximal ideal $tB'$, and is
Noetherian because it is finite over $B$, therefore a DVR. This contradicts the
maximality of $B$, so that we must have $\varphi(B) = K$.
\end{proof}

\begin{remark}
Since $B$ is an integral domain containing $A$, it is flat over $A$ by
Ex.~\ref{ex:10.2-ex}. In EGA $0_{IV}$ (10.3.1), the following more general fact is proved:
let $(A, m, k)$ be a Noetherian local ring, and $K$ an extension field of $k$, then
there exists a Noetherian ring $B$ containing $A$ satisfying the three conditions
(1) $\rad(B) = mB$, (2) $B/mB$ is isomorphic over $k$ to $K$, and (3) $B$ is flat over $A$.
\end{remark}

\begin{theorem}\label{thm:29.2}
Let $(A, m, K)$ be a complete local ring, $(R, pR, k)$ a $p$-ring, and
$\varphi_0\colon k \to K$ a field homomorphism; then there exists a local homomorphism
$\varphi\colon R \to A$ which induces $\varphi_0$ on the residue fields.
\end{theorem}

\begin{proof}
Set $S = \mathbb{Z}_{(p)}$, and let $k_0 \subset k$ be the prime subfield. Since
$\varphi_0(k_0) \subset K$, the prime number $p$, viewed as an element of $A$, belongs to
$m$. Hence the standard homomorphism $\mathbb{Z} \to A$ extends to a local
homomorphism $S \to A$. Now $R \otimes_S k_0 = R/pR = k$ is a separable extension of
$k_0$, and hence 0-smooth over $k_0$; also $R$ is a torsion-free $S$-module, hence
flat over $S$, so that by Theorem~\hyperref[thm:28.10]{28.10}, $R$ is $pR$-smooth over $S$.
\[
\begin{tikzcd}
k_0 \arrow[r, "\subset"] \arrow[d, "\subset"'] & k \arrow[r, "\varphi_0"] & K = A/m \\
S \arrow[r, "\subset"] & R \arrow[ru, "\varphi"', dashed] & A \arrow[u, "p"] \\
& & \mathbb{Z} \arrow[u, "\subset"] \arrow[uu, "\subset"]
\end{tikzcd}
\]
Hence, as we discussed at the beginning of \S\ref{sec:28}, we can lift
$\varphi_0\colon k \to A/m = K$ successively to $R \to A/m^\nu$, and using the fact that
$A = \varprojlim A/m^\nu$, we get $\varphi\colon R \to A$ making the left-hand diagram
commute.
\end{proof}

\begin{theorem}\label{thm:29.3}
The complete $p$-ring is uniquely determined up to isomorphism by its residue
field.
\end{theorem}

\begin{proof}
Suppose that $R$ and $R'$ are both complete $p$-rings with residue field $k$; then by
the theorem there exists a local homomorphism $\varphi\colon R \to R'$ which induces
the identity map on the residue field. We have $R' = \varphi(R) + pR'$, and of course
$\varphi(p) = p$, so that by the completeness of $R$ we see that $\varphi$ is
surjective, and is also injective, since $p^\nu R$ is not contained in $\Ker \varphi$
for any $\nu$. Therefore $R \cong R'$.
\end{proof}

Let $(A, m, k)$ be a complete local ring of unequal characteristic, and let
$p = \chark k$. We say that a subring $A_0 \subset A$ is a \emph{coefficient ring}\index{coefficient ring}
of $A$ if $A_0$ is a complete Noetherian local ring with maximal ideal $pA_0$, and
\[
  A = A_0 + m, \quad \text{that is, } k = A/m \cong A_0/pA_0.
\]
By Theorem~\hyperref[thm:29.1]{1} applied to the residue field $k$ of $A$, there exists a $p$-ring
$S$ such that $S/pS = k$; write $R$ for the completion of $S$, so that $R$ is a
complete $p$-ring with residue field $k$. By Theorem~\hyperref[thm:29.2]{2}, there exists a local
homomorphism $\varphi\colon R \to A$ inducing an isomorphism on the residue fields. If
we set $\varphi(R) = A_0$, then this is clearly a coefficient ring of $A$. If $A$ has
characteristic $0$ then $\varphi$ is injective and $A_0 \cong R$. If $A$ has
characteristic $p^\nu$ then $A_0 \cong R/p^\nu R$. We summarise the above discussion in
the following theorem.

\begin{theorem}\label{thm:29.4}
If $(A, m, k)$ is a complete local ring and $p = \chark k$ then $A$ has a coefficient
ring $A_0$. If $A$ has characteristic $0$ then $A_0$ is a complete DVR, and if $A$ has
characteristic $p^\nu$ then $A_0$ is an Artinian local ring.
\end{theorem}

In what follows, in order to include the case of equal characteristic in our
discussion, we also consider coefficient fields as being coefficient rings.
From the previous theorem and Theorem~\hyperref[thm:28.3]{28.3}, we get the following
important result.

\begin{theorem}\label{thm:29.5}
\begin{enumerate}[(i)]
  \item If $(A, m)$ is a complete local ring and $m$ is finitely generated, then $A$
        is Noetherian.
  \item A Noetherian complete local ring is a quotient of a regular local ring;
        in particular it is universally catenary.
  \item If $A$ is a Noetherian complete local ring (and in the case of unequal
        characteristic, $A$ is an integral domain), then there exists a subring
        $A' \subset A$ with the following properties: $A'$ is a complete regular local
        ring with the same residue field as $A$, and $A$ is finitely generated as an
        $A'$-module.
\end{enumerate}
\end{theorem}

\begin{proof}
We choose a coefficient ring $A_0$ of $A$. If $m = (x_1, \ldots, x_n)$ then every element
of $A$ can be expanded as a power series in $x_1, \ldots, x_n$ with coefficients in
$A_0$, so that $A$ is a quotient of $A_0[[X_1, \ldots, X_n]]$ and hence Noetherian. Now
$A_0$ is a quotient of a $p$-ring $R$, so that $A$ is a quotient of
$R[[X_1, \ldots, X_n]]$, which is a regular local ring, hence a CM ring, and therefore
according to Theorem~\hyperref[thm:17.9]{17.9}, $A$ is universally catenary. To prove (iii), set
$\dim A = n$, and in the case of equal characteristic let $\{y_1, \ldots, y_n\}$ be any
system of parameters of $A$; if $A$ is an integral domain of characteristic $0$ and
$\chark k = p$, we can choose a system of parameters $\{y_1 = p, y_2, \ldots, y_n\}$ of $A$
starting with $p$. In either case $R = A_0$, so that we set $A' = R[y]$. Then $A'$ is the
image of $\varphi\colon R[Y] \to A$, where $R[Y]$ is the regular local ring
\[
  R[Y] = R[[Y_1, \ldots, Y_n]], \quad \text{or } R[[Y_2, \ldots, Y_n]] \text{ if } y_1 = p,
\]
and $\varphi$ is the $R$-algebra homomorphism defined by $\varphi(Y_i) = y_i$. Set
$m' = \sum y_iA'$. Since $A/m = A'/m'$, every $A$-module of finite length has the same
length when viewed as an $A'$-module. In particular $A/m'A$ is a finite module over
$A'/m'$, and $A$ is $m'$-adically separated, so that by Theorem~\hyperref[thm:8.4]{8.4} $A$ is a finite
$A'$-module. Therefore, we have
\[
  \dim A' = \dim A = n.
\]
$R[Y]$ is an $n$-dimensional integral domain, and if $\Ker \varphi \ne 0$ we would have
$\dim A' < n$, which is a contradiction. Therefore $\varphi$ is injective and
$A' \cong R[Y]$.
\end{proof}

\begin{remark}
In the case of unequal characteristic when $A$ is not an integral domain, (iii)
can fail to hold. If $A$ is of characteristic $p^m$ with $m > 1$, then every subring
of $A$ has the same characteristic $p^m$, so cannot be regular. Even if $A$ has
characteristic $0$, the following is a counter-example: let $R$ be a complete
$p$-ring and $A = R[X]/(pX)$; then $A$ is a complete one-dimensional Noetherian
local ring, but if a subring $A'$ as in (iii) were to exist, $A'$ would be a
one-dimensional regular local ring, hence a DVR, and since $A'$ has characteristic
$0$ and its residue field characteristic $p$, $A'/pA'$ would be an Artinian ring,
and hence also $A/pA$ would be Artinian. But $A/pA \cong k[X]$ is
one-dimensional, and this is a contradiction.

The proof of Theorem~\hyperref[thm:29.5]{5} shows that it is sufficient to assume that $p$ is not
in any minimal prime ideal of $A$.
\end{remark}

Corresponding to the definition of quasi-coefficient field of an
equicharacteristic local ring, let us define \emph{quasi-coefficient rings} in the case of
unequal characteristic. Let $(A, m, K)$ be a possibly non-complete local ring, and
suppose $\chark K = p$. A subring $S \subset A$ is said to be a quasi-coefficient ring if
it satisfies the following two conditions:
\begin{enumerate}[(1)]
  \item $S$ is a Noetherian local ring with maximal ideal $pS$;
  \item $K = A/m$ is 0-\'etale over $S/pS$.
\end{enumerate}
In view of (1) if $A$ has characteristic $0$ then $S$ is a DVR, and if $A$ has
characteristic $p^\nu$ then $S$ is an Artinian ring.

\begin{theorem}\label{thm:29.6}
Let $(A, m, K)$ be a local ring, and suppose $\chark K = p$. Let $C \subset A$ be a
subring, and assume that $C$ is a Noetherian local ring with maximal ideal $pC$,
and that $K = A/m$ is separable over $C/pC$. Then there exists a quasi-coefficient
ring $S$ of $A$ containing $C$; moreover, if $A$ is flat over $C$, then it is also
flat over $S$.
\end{theorem}

\begin{proof}
Let $\{B_\lambda\}_{\lambda \in \Lambda}$ be a $p$-basis of $K$ over $C/pC$, and choose an inverse
image $b_\lambda \in A$ for each $\beta_\lambda$. Setting $C[\{b_\lambda\}] = C'$, by
Theorem~\hyperref[thm:26.8]{26.8} we see that $C'/(m \cap C') = (C/pC)[\{\bar{X}_\lambda\}]$ is a
polynomial ring over $C/pC$. Hence if $f(\ldots, X_\lambda, \ldots)$ is a non-zero
polynomial with coefficients in $C$ which satisfies $f(b_\lambda) \in m$, then setting
$\rho$ for the highest common factor of the coefficients of $f$, we have
\[
  f(X) = \rho \cdot f_1(X) \quad \text{and} \quad f_1(b_\lambda) \not\equiv 0.
\]
Thus $f_1(b_\lambda) \notin m$, and we must have $\rho \in m$, so that we have shown that
$m \cap C' = pC'$. Setting $S = (C')_{m}$, we have $S \subset A$, $m \cap S = pS$ and
$S/pS = K'$; since $C'$ is $p$-adically separated, so is $S$, and hence all the
ideals of $S$ are of the form $(0)$ or $(p^\nu)$. Thus $S$ is Noetherian, and it
satisfies all the conditions for a quasi-coefficient ring of $A$. If $A$ is flat over
$C$, then
\[
  p^\nu C \otimes_C A \cong p^\nu A.
\]
and hence the composite
\[
  p^\nu C \otimes A = (p^\nu C \otimes_C S) \otimes_S A \to p^\nu S \otimes_S A \to p^\nu A
\]
is injective; but the first arrow is surjective, so that the second arrow
$p^\nu S \otimes_S A \to p^\nu A$ is injective. By Theorem~\hyperref[thm:7.7]{7.7}, this proves that $A$ is
flat over $S$.
\end{proof}

All the assumptions of this theorem are satisfied by taking $C$ to be the image
of $\mathbb{Z}_{(p)} \to A$, so that this proves that every local ring has a quasi-coefficient
ring (including the quasi-coefficient field of a local ring in the equal
characteristic case).

\begin{theorem}\label{thm:29.7}
Let $(A, m, K)$ be a local ring, and $\hat{A}$ its completion, and suppose
$\chark K = p$. Let $S$ be a quasi-coefficient ring of $A$, and write $\bar{S}$ for its
image in $\hat{A}$; then there exists a unique coefficient ring $A_0$ of $\hat{A}$
containing $\bar{S}$.
\end{theorem}

\begin{proof}
Since $\bar{S}$ is a quasi-coefficient ring of $\hat{A}$, we can assume that $A$ is a
complete local ring. If $A$ has characteristic $0$ then $S$ is a DVR, and by
Theorem~\hyperref[thm:29.1]{1} there exists a complete $p$-ring $R$ containing $S$ and with
residue field $K$. Now $K$ is 0-\'etale over $S/pS$, and $R$ is flat over $S$, so that
$R$ is $pR$-\'etale over $S$ (by Ex.~\ref{ex:28.1-ex}). Hence there exists a unique
$S$-algebra homomorphism $R \to A$ which induces the identity map on the residue
fields; we write $A_0$ for the image of $R$. Then $R \cong A_0$, and $A_0$ is a
coefficient ring of $A$. If $A$ had another coefficient ring $B$ containing $S$ then
$B$ would also be a complete $p$-ring, and for the same reason as above, there
exist unique $S$-algebra homomorphisms $B \to A_0$ and $B \to A_0$, so that we must
have $B = A_0$.

If $A$ has characteristic $p^\nu$ then $S$ is an Artinian ring, and therefore
complete, so that applying Theorem~\hyperref[thm:29.4]{4} to $S$ itself, we can write
$S = R_0/p^\nu R_0$ with $R_0$ a complete $p$-ring. Let $R$ be a complete $p$-ring
containing $R_0$ and with residue field $K$; then by Ex.~\ref{ex:28.1-ex}, $R$ is $pR$-\'etale
over $R_0$, so that there exists a unique $R_0$-algebra homomorphism $R \to A$
inducing the identity map on the residue field $K$. The image is a coefficient
ring of $A$ containing $S$; the uniqueness is proved as in the case of
characteristic $0$.
\end{proof}

Next we study the structure of complete regular local rings. Let $(A, m, K)$ be a
local ring of unequal characteristic, and suppose that $\chark K = p$; then $A$ is
said to be \emph{ramified}\index{ramified} if $p \in m^2$ and \emph{unramified} if $p \notin m^2$. We will
also say that $A$ is unramified in the case of equal characteristic.

\begin{theorem}\label{thm:29.8}
An unramified complete regular local ring is a formal power series ring over a
field or over a complete $p$-ring.
\end{theorem}

\begin{proof}
Let $R$ be a coefficient ring of $A$. In the case of equal characteristic, $R$ is a
field, and if $x_1, \ldots, x_n$ is a regular system of parameters of $A$ then
$A = R[x_1, \ldots, x_n] \cong R[[X_1, \ldots, X_n]]$ (see the proof of
Theorem~\hyperref[thm:29.5]{5}). In the case of unequal characteristic, $R$ is a complete $p$-ring,
and since $p \notin m^2$, we can choose a regular system of parameters
$\{p, x_2, \ldots, x_n\}$ of $A$ containing $p$. Then $A = R[x_2, \ldots, x_n] \cong
R[[X_2, \ldots, X_n]]$.
\end{proof}

In the ramified case, it is not necessarily the case that $A$ can be expressed as
a formal power series ring over a DVR. To give the structure theorem in this
case we need the notion of an \emph{Eisenstein extension}.

\begin{lemma}[\textbf{Eisenstein's irreducibility criterion}]\label{lem:29.1}
Let $A$ be a ring, and $f(X) = X^n + a_1X^{n-1} + \cdots + a_n$, with $a_i \in A$. If there
exists a prime ideal $P$ of $A$ such that $a_1, \ldots, a_n \in P$ but $a_n \notin P^2$ then $f$
is irreducible in $A[X]$. If in addition $A$ is an integrally closed domain then
the principal ideal $(f)$ is a prime ideal of $A[X]$.
\end{lemma}

\begin{proof}
If $f$ is reducible, we can write
\[
  f = (X^r + b_1X^{r-1} + \cdots + b_r)(X^s + c_1X^{s-1} + \cdots + c_s)
\]
with $0 < r < n$, $s = n - r$ and $b_i, c_j \in A$. Reducing the coefficients on either
side modulo $P$ we have
\[
  X^n = (X^r + \bar{b}_1X^{r-1} + \cdots + \bar{b}_r)(X^s + \bar{c}_1X^{s-1} + \cdots + \bar{c}_s)
\]
in $(A/P)[X]$, so that we must have $b_i, c_j \in P$ for all $i, j$, but then
$a_n = b_rc_s \in P^2$, which contradicts the assumption. If $A$ is an integrally closed
domain and $K$ is the field of fractions of $A$, then by Ex.~\ref{ex:9.6-ex}, $f$ remains
irreducible in $K[X]$. Also, $f$ is manic, so that we have
$f \cdot A[X] = f \cdot K[X] \cap A[X]$, and this is a prime ideal of $A[X]$.
\end{proof}

Let $(A, m)$ be a normal local ring; then an extension ring
\[
  B = A[x]/(f) = A[x]
\]
defined by an \emph{Eisenstein polynomial}\index{Eisenstein polynomial}
\[
  f = X^n + a_1X^{n-1} + \cdots + a_n \quad \text{with } a_i \in m \text{ for all } i,
      \text{ and } a_n \notin m^2
\]
is called an \emph{Eisenstein extension} of $A$. By the lemma, $B$ is an integral domain,
and is integral over $A$. We have $B/mB = (A/m)[X]/(X^n)$, so that $B$ has just one
maximal ideal $n = mB + xB$. Hence $B$ is a local ring, and its residue field
coincides with that of $A$.

\begin{theorem}\label{thm:29.9}
\begin{enumerate}[(i)]
  \item If $(A, m)$ is a regular local ring, then an Eisenstein extension of $A$ is
        again a regular local ring.
  \item If $A$ is a ramified complete regular local ring and $R$ is a coefficient
        ring of $A$ then there is a subring $A_1 \subset A$ with the following
        properties:
        \begin{enumerate}[(1)]
          \item $A_1$ is an unramified complete regular local ring containing $R$,
                and hence can be expressed as a formal power series ring over $R$;
          \item $A$ is an Eisenstein extension of $A_1$.
        \end{enumerate}
\end{enumerate}
\end{theorem}

\begin{proof}
(i) Let $B = A[x]$, and $x^n + a_1x^{n-1} + \cdots + a_n = 0$, with $a_i \in m$ and
$a_n \in m^2$ but $a_n \notin m^3$. Then there exists a regular system of parameters
$\{y_1, \ldots, y_n = a_n\}$ of $A$ with $a_n$ as an element. As we have seen above, the
maximal ideal of $B$ is $mB + xB$, but $a_n \in xB$, so that
$\{y_1, \ldots, y_{n-1}, x\}$ is a regular system of parameters of $B$.

(ii) Since $\htop pA = 1$, by a skilful choice of a regular system of parameters
$\{x_1, \ldots, x_d\}$ of $A$, we can arrange that $\{p, x_2, \ldots, x_d\}$ is a system of
parameters of $A$. If we set $A_1 = R[x_2, \ldots, x_d]$ then $A_1$ is a complete
unramified regular local ring, and $A$ is a finite module over $A_1$ (see the
proof of Theorem~\hyperref[thm:29.5]{5}). We set $m_1$ for the maximal ideal of $A_1$. Now
$A = m_1A + A_1[x_1]$, so that by Theorem~\hyperref[thm:8.4]{8.4} (or by NAK),
$A = A_1[x_1]$. Let
\[
  f(X) = X^n + a_1X^{n-1} + \cdots + a_n \quad \text{with } a_i \in A_1
\]
be the minimal polynomial of $x_1$ over $A_1$. Then $a_n = x_1 \cdot \ldots \in m$,
so that $a_n \notin m_1^2$. Therefore by Hensel's lemma (Theorem~\hyperref[thm:8.3]{8.3}), all the
$a_i \in m_1$. We are left to prove that $a_n \notin m_1^2$. Write $p = \sum_{i=1}^{d} b_ix_i$ with
$b_i \in A_1$, and express the $b_i$ in the form $b_i = q_i(x_1)$, with
$q_i(X) \in A_1[X]$; then $x_1$ is a root of
\[
  F(X) = q_1(X)X + \sum_{i=2}^{d} q_i(X)x_i - p
\]
so that $F(X)$ is divisible by $f(X)$. Hence the constant term $F(0)$ of $F$ is
divisible by $a_n$. However, $F(0) = \sum_{i=2}^{d} q_i(0)x_i - p$, and
$p, x_2, \ldots, x_d$ is a regular system of parameters, so that $F(0) \notin m^2$, hence
also $a_n \notin m_1^2$.
\end{proof}

% ============ 习题 ============
\section*{Exercises to \S\ref{sec:29}}
\emph{Prove the following propositions.}

\begin{exercise}\label{ex:29.1-ex}
Let $A$ be a complete $p$-ring, $y$ an indeterminate over $A$, and $B = A[[y]]$; let
$C = B[x]$ be the Eisenstein extension of $B$ given by $x^2 + yx + p = 0$. Then $C$
is a two-dimensional complete regular local ring, but is not a formal power
series ring over a DVR of characteristic $0$.
\end{exercise}

\begin{exercise}\label{ex:29.2-ex}
In Theorem~\hyperref[thm:29.2]{29.2}, if $k$ is a perfect field then $\varphi$ is uniquely determined
by $\varphi_0$.
\end{exercise}
